% !TEX root = ../thesis.tex
% =============================================================================
% Chapter 3: Behavioural Feature Design
% =============================================================================

\chapter{Behavioural Feature Design}
\label{ch:features}

This chapter defines the behavioural feature set used throughout the thesis. \Cref{sec:features:principles} states the criteria that guided feature design and selection. The remaining sections present the 32 resulting features, organised by origin and function: inherited metrics, step-size and movement dynamics, information-theoretic measures, adapted population dynamics, and features new to the optimisation literature.

Every feature is computed from a single trajectory, the ordered sequence of candidate solutions that an algorithm evaluates together with their fitness values. Formally, a trajectory is a sequence $\{(x_t, y_t)\}_{t=1}^{T}$ where $x_t \in \mathbb{R}^d$ is the $t$-th candidate solution with decision variables $x_t^{(0)}, \ldots, x_t^{(d-1)}$, and $y_t = \mathcal{F}(x_t)$ is its fitness. The best-so-far fitness at evaluation~$t$ is denoted $f^*_t = \min_{s \le t} y_s$. Throughout this thesis, trajectories are collected over $T = 2{,}000 \times d$ evaluations in $d = 5$ dimensions with bounds $[-5, 5]$, giving $T = 10{,}000$ evaluations per trajectory.

% ============================================================================
\section{Design Principles}
\label{sec:features:principles}

Four criteria guided the selection and design of behavioural features.

\begin{enumerate}
  \item \textbf{Computational cheapness.} MA-BBOB evaluations already carry the bulk of experiment runtime. Behavioural features should add only a small fraction of additional expense.

  \item \textbf{Distinctness.} Each feature should capture a dimension of behaviour not already covered by other features.

  \item \textbf{Interpretability for LLM feedback.} Features are reported to the language model as part of the behavioural feedback string. They must therefore carry clear, concise semantics that an LLM can act on.

  \item \textbf{Discrimination.} Features should separate high-performing from low-performing algorithms. We quantify this via Spearman correlation ~\cite{spearman1904proof} with AOCC and Kolmogorov--Smirnov effect size ~\cite{kolmogorov1933} between the top-25\% and bottom-25\% candidates.
\end{enumerate}

% ============================================================================
\section{Existing BLADE Metrics}
\label{sec:features:blade}

BLADE~\cite{vanStein2025BLADE} ships with eleven behavioural metrics introduced by van~Stein et al.~\cite{vanStein2025Behaviour}. They fall into four functional categories covering exploration and diversity, exploitation and intensification, convergence, and stagnation. We define each metric formally below, using the trajectory notation from the start of the chapter.

\paragraph{\texttt{avg\_nearest\_neighbor\_distance}.}
For each evaluation $x_k$ (sampled every tenth step for efficiency), the Euclidean distance to its nearest neighbour among the $H = 1{,}000$ most recent prior evaluations:
\begin{equation}
\label{eq:avg_nn_dist}
  \text{ANN} = \frac{1}{|\mathcal{K}|} \sum_{k \in \mathcal{K}} \min_{\max(0,\, k-H) \le j < k} \lVert x_k - x_j \rVert,
\end{equation}
where $\mathcal{K} = \{1, 11, 21, \ldots\}$. Large values indicate the search covers new territory at each step. Small values indicate repeated returns to previously visited regions.

\paragraph{\texttt{dispersion}.}
The radius of the largest empty hypersphere that fits inside $\mathcal{S}$ without containing any evaluated point:
\begin{equation}
\label{eq:dispersion_formal}
  \mathrm{disp}(X) = \sup_{y \in \mathcal{S}} \min_{1 \le t \le T} \lVert y - x_t \rVert_2.
\end{equation}
In practice the supremum is approximated by drawing $N = 10{,}000$ points $z_i$ uniformly from $\mathcal{S}$ and computing
\begin{equation}
\label{eq:dispersion}
  \widehat{\mathrm{disp}}(X) = \max_{1 \le i \le N} \min_{1 \le t \le T} \lVert z_i - x_t \rVert_2,
\end{equation}
with nearest-neighbour queries accelerated by a KD-tree. Large values indicate that large regions of $\mathcal{S}$ remain unvisited. Lower values signal better search-space coverage.

\paragraph{\texttt{avg\_exploration\_pct}.}
For any finite set $S \subset \mathbb{R}^d$, the average pairwise Euclidean distance is
\begin{equation}
\label{eq:pairwise_dist}
  D(S) = \frac{2}{|S|(|S|-1)} \sum_{\substack{p < q \\ x_p, x_q \in S}} \lVert x_p - x_q \rVert_2.
\end{equation}
With axis widths $w_k = u_k - l_k$ and mean width $\bar{w} = \frac{1}{d}\sum_{k=1}^d w_k$, the expected pairwise distance of two i.i.d.\ uniform samples in $\mathcal{S}$ is approximated by
\begin{equation}
\label{eq:rand_div}
  D_{\mathrm{rand}} = \bar{w} \sqrt{d/6}.
\end{equation}
The trace $X$ is partitioned into $C$ consecutive chunks $S_0, \ldots, S_{C-1}$ of size $K = 500$, and each chunk is scored relative to random search:
\begin{equation}
\label{eq:chunk_exploration}
  E_c = \min\!\left\{100 \cdot \frac{D(S_c)}{D_{\mathrm{rand}}},\, 100\right\} \quad (0 \le E_c \le 100).
\end{equation}
The average exploration percentage is the mean across chunks:
\begin{equation}
\label{eq:avg_expl_pct}
  \bar{E} = \frac{1}{C} \sum_{c=0}^{C-1} E_c.
\end{equation}
Values near 100 indicate chunk-level diversity on par with random search. Values near 0 indicate concentrated, exploitative search.

\paragraph{\texttt{avg\_distance\_to\_best}.}
Let $b^*(k) = \arg\min_{t \le k} y_t$ denote the index of the best-so-far point at evaluation $k$. The mean distance from each evaluation to the best-so-far incumbent:
\begin{equation}
\label{eq:avg_dist_best}
  \bar{\Delta} = \frac{1}{T-1} \sum_{k=1}^{T-1} \lVert x_k - x_{b^*(k)} \rVert.
\end{equation}
Low values indicate strong local exploitation around the incumbent.

\paragraph{\texttt{intensification\_ratio}.}
The fraction of evaluations within a radius $r = 0.1 \cdot \bar{w}$ of the final best point $x^* = \arg\min_{t} y_t$:
\begin{equation}
\label{eq:intensification}
  I = \frac{1}{T} \sum_{t=1}^{T} \mathbf{1}\!\left[\lVert x_t - x^* \rVert < r\right].
\end{equation}
High values indicate strong exploitation near the eventual best solution.

\paragraph{\texttt{avg\_exploitation\_pct}.}
The complement of the exploration percentage:
\begin{equation}
\label{eq:avg_exploit_pct}
  \bar{X} = 100 - \bar{E}.
\end{equation}
It captures the share of chunk-level variance attributable to local refinement.

\paragraph{\texttt{average\_convergence\_rate}.}
Let $e_t = f^*_t - f^*_T$ be the error of the best-so-far above the final best, and $\varepsilon$ machine epsilon (added to avoid division by zero). The geometric mean of successive error ratios:
\begin{equation}
\label{eq:acr}
  \text{ACR} = \exp\!\left( \frac{1}{T-1} \sum_{t=1}^{T-1} \log \frac{e_{t+1} + \varepsilon}{e_t + \varepsilon} \right).
\end{equation}
$\text{ACR} < 1$ indicates convergence. Smaller values indicate faster convergence.

\paragraph{\texttt{avg\_improvement}.}
Let $\mathcal{I} = \{t : y_t < f^*_{t-1}\}$ be the set of improving iterations. The mean improvement magnitude over those iterations:
\begin{equation}
\label{eq:avg_imp}
  \bar{I} = \frac{1}{|\mathcal{I}|} \sum_{t \in \mathcal{I}} \bigl(f^*_{t-1} - y_t\bigr).
\end{equation}
Large values indicate the algorithm makes substantial gains when it improves. Small values indicate refinement in small increments.

\paragraph{\texttt{success\_rate}.}
The fraction of iterations that improve the best-so-far:
\begin{equation}
\label{eq:success_rate}
  S = \frac{|\mathcal{I}|}{T - 1}.
\end{equation}
High values indicate the algorithm frequently finds better solutions, which may reflect either an easy problem or effective search steps. Low values indicate the algorithm struggles to improve on the incumbent.

\paragraph{\texttt{longest\_no\_improvement\_streak}.}
The length of the longest contiguous run of iterations that do not improve the best-so-far:
\begin{equation}
\label{eq:longest_no_imp}
  L = \max \bigl\{b - a + 1 : y_t \ge f^*_{t-1} \text{ for all } t \in [a, b]\bigr\}.
\end{equation}
Streaks of several thousand out of $T = 10{,}000$ indicate the algorithm has effectively stopped progressing.

\paragraph{\texttt{last\_improvement\_fraction}.}
Let $t^* = \max \mathcal{I}$ be the index of the last improving iteration. The fraction of the budget that elapses after it:
\begin{equation}
\label{eq:last_imp_frac}
  \Phi = \frac{T - 1 - t^*}{T - 1}.
\end{equation}
High values indicate the algorithm plateaued early and made no further progress. Low values indicate it continued improving until near the end.

% ============================================================================
\section{Step-Size and Movement Dynamics}
\label{sec:features:stepsize}

The first group of new features characterises the basic mechanics of how
a candidate algorithm moves through decision space.  Step-size statistics
(mean, standard deviation, OLS slope) and directional persistence are
basic trajectory descriptors.  Step-size statistics in particular are
common in time-series feature libraries such as
\texttt{tsfresh}~\cite{christ2018tsfresh}. De~Nobel et
al.~\cite{deNobel2021tsfresh} applied \texttt{tsfresh} to CMA-ES internal
strategy parameters (step-size $\sigma$, evolution paths), though
notably none of the simple statistics over $\sigma$ that they considered
were retained by their feature selector.  We compute analogous statistics
directly on Euclidean step magnitudes $\lVert x_{t+1} - x_t \rVert$ from
the trajectory, which is algorithm-agnostic rather than CMA-ES-specific.
Directional persistence, the mean cosine similarity between consecutive
displacement vectors, is a standard descriptor in random-walk and
movement-ecology analyses~\cite{codling2008randomwalk}.  Let
$s_t = \lVert x_{t+1} - x_t \rVert$ denote the Euclidean step size at
evaluation~$t$.

\paragraph{\texttt{step\_size\_mean}.}
The mean step size across the trajectory:
\begin{equation}
\label{eq:step_mean}
  \bar{s} = \frac{1}{T-1}\sum_{t=1}^{T-1} s_t.
\end{equation}
Large values of the mean step size~$\bar{s}$ indicate broad exploration.
Small values indicate local exploitation.

\paragraph{\texttt{step\_size\_std}.}
The standard deviation of step sizes:
\begin{equation}
\label{eq:step_std}
  \sigma_s = \sqrt{\frac{1}{T-2}\sum_{t=1}^{T-1}
    \bigl(s_t - \bar{s}\bigr)^2}.
\end{equation}
High step-size standard deviation~$\sigma_s$ signals alternating
exploration and exploitation phases. Low variability signals a uniform
step-size regime throughout the run.

\paragraph{\texttt{step\_size\_trend}.}
The ordinary least-squares slope of step size against evaluation index:
\begin{equation}
\label{eq:step_trend}
  \beta_s
    = \frac{\sum_{t}(t - \bar{t})(s_t - \bar{s})}
           {\sum_{t}(t - \bar{t})^2}.
\end{equation}
A negative step-size trend~$\beta_s$ indicates the classic convergence
pattern of contracting step sizes. A positive slope suggests the
algorithm diverges over time.

\paragraph{\texttt{directional\_persistence}.}
The mean cosine similarity between consecutive displacement vectors
$d_t = x_{t+1} - x_t$:
\begin{equation}
\label{eq:dir_persist}
  \phi = \frac{1}{T-2}\sum_{t=2}^{T-1}
    \frac{d_t \cdot d_{t-1}}
         {\lVert d_t \rVert \, \lVert d_{t-1} \rVert}.
\end{equation}
Values of the directional persistence~$\phi$ near~$+1$ indicate that the
search maintains a consistent heading
(gradient-following behaviour). Values near~$0$ indicate a random walk.
Negative values indicate anti-persistent zigzag motion.

% ============================================================================
\section{Information-Theoretic Features}
\label{sec:features:infotheo}

The four features in this category apply established time-series complexity
measures to the raw fitness sequence $\{y_1, \ldots, y_T\}$.  While each
measure is well-known in signal processing and physics, their application to
optimisation algorithm fitness trajectories on BBOB benchmarks is new.

\paragraph{\texttt{fitness\_sample\_entropy}.}
Sample Entropy (SampEn) quantifies the regularity of the fitness
series~\cite{richman2000sampen}.  Given embedding dimension~$m$ and tolerance
$r = 0.2 \cdot \mathrm{std}(y)$, it counts the conditional probability that
sequences matching for $m$~consecutive values still match at length $m+1$:
\begin{equation}
\label{eq:sampen}
  \mathrm{SampEn}(m, r)
    = -\ln \frac{A^m(r)}{B^m(r)},
\end{equation}
where $B^m$ is the number of template matches of length~$m$ and $A^m$ the
number that also match at length $m+1$ (excluding self-matches).  Low SampEn
indicates a repetitive, self-similar convergence pattern. High SampEn indicates
complex, unpredictable dynamics.  We use $m = 2$ and subsample every 10th
evaluation ($N = 1{,}000$) to reduce the $O(N^2 m)$ cost from ${\sim}30$\,s to
${\sim}0.3$\,s per trajectory.

\paragraph{\texttt{fitness\_permutation\_entropy}.}
Normalised Permutation Entropy (PE) maps the fitness series to a sequence of
ordinal patterns and measures their distributional uniformity~\cite{bandt2002permutation}.
Each window of $D$~consecutive values is replaced by its rank permutation
$\pi \in S_D$, and PE is the Shannon entropy of the permutation distribution
normalised by $\log_2(D!)$:
\begin{equation}
\label{eq:pe}
  \mathrm{PE}(D)
    = -\frac{1}{\log_2(D!)}\sum_{\pi \in S_D} p(\pi) \log_2 p(\pi).
\end{equation}
We use $D = 5$, giving $D! = 120$ possible patterns against $T = 10{,}000$
data points.  $\mathrm{PE} \approx 0$ corresponds to perfectly monotone
convergence. $\mathrm{PE} \approx 1$ to a completely random sequence.

\paragraph{\texttt{fitness\_autocorrelation}.}
The lag-1 autocorrelation of the raw fitness series, following Weinberger's
framework for fitness landscape correlation structure~\cite{weinberger1990fitness}:
\begin{equation}
\label{eq:fitac}
  \rho_y(1)
    = \frac{\sum_{t=1}^{T-1}(y_t - \bar{y})(y_{t+1} - \bar{y})}
           {\sum_{t=1}^{T}(y_t - \bar{y})^2}.
\end{equation}
Importantly, while Weinberger's autocorrelation has long been used to
characterise landscapes via \emph{separate} random walks, it has not previously
been computed on the algorithm's \emph{own} trajectory on BBOB.  High positive
$\rho_y(1)$ indicates smooth, locally exploitative search. Values near zero
indicate the algorithm jumps between unrelated fitness regions.

\paragraph{\texttt{fitness\_lempel\_ziv\_complexity}.}
Normalised Lempel--Ziv complexity (LZC) measures the rate at which new
patterns appear in a binarised version of the fitness
series~\cite{lempel1976complexity}.  The binarisation maps each step to
$b_t = \mathbf{1}[y_{t+1} < y_t]$ (improvement or not), yielding a binary
string.  LZ76 parsing counts the number of distinct substrings~$c(N)$, and the
normalised complexity is:
\begin{equation}
\label{eq:lzc}
  \mathrm{LZC}_n = \frac{c(N)}{N / \log_2 N}.
\end{equation}
$\mathrm{LZC}_n \approx 1$ indicates a random improvement/stagnation pattern.
$\mathrm{LZC}_n \approx 0$ indicates highly repetitive structure (e.g.,
long unbroken stagnation followed by a burst of improvements).

% ============================================================================
\section{Adapted Population Dynamics}
\label{sec:features:dynamo}

Cenikj et al.~\cite{cenikj2023dynamorep} introduced DynamoRep, a set of
trajectory-based features that characterise how a population's spatial and
fitness distributions change across generations.  Their features assume a
population-based algorithm with multiple individuals per generation.  In our
$(1{+}1)$-ES setting the population size is~1, so per-generation statistics
are undefined.  We adapt the DynamoRep idea by comparing the first 25\% of
evaluations (\emph{early}) with the last 25\% (\emph{late}), treating each
window as a proxy for the population at that stage of the search.

Let $X_{\mathrm{early}} = \{x_t : t \le T/4\}$ and
$X_{\mathrm{late}} = \{x_t : t > 3T/4\}$.

\paragraph{\texttt{x\_spread\_early} and \texttt{x\_spread\_late}.}
Mean per-dimension standard deviation of decision variables in the early and
late windows, respectively:
\begin{equation}
\label{eq:spread}
  S_{\mathrm{E}}
    = \frac{1}{d}\sum_{j=0}^{d-1}
      \mathrm{std}\!\bigl(\{x_t^{(j)} : t \le T/4\}\bigr).
\end{equation}
The late spread~$S_{\mathrm{L}}$ is defined identically over $t > 3T/4$.
High early spread~$S_{\mathrm{E}}$ indicates broad initial exploration.
Low~$S_{\mathrm{L}}$ indicates the search has converged to a small
region.

\paragraph{\texttt{spread\_ratio}.}
The ratio of late to early spread:
\begin{equation}
\label{eq:spread_ratio}
  R_S = S_{\mathrm{L}} \;/\; S_{\mathrm{E}}.
\end{equation}
Values of the spread ratio~$R_S$ below~1 indicate the algorithm narrows
its search over time (the expected convergence pattern). Values above~1
indicate late-stage diversification.

\paragraph{\texttt{centroid\_drift}.}
The Euclidean distance between the centroids of the early and late windows:
\begin{equation}
\label{eq:centroid}
  D_c
    = \bigl\lVert \bar{x}_{\mathrm{early}}
      - \bar{x}_{\mathrm{late}} \bigr\rVert.
\end{equation}
Large centroid drift~$D_c$ means the algorithm relocated its search
centre during the run it discovered a new promising region.  Small drift means the algorithm stayed
local throughout.

\paragraph{\texttt{f\_range\_early}, \texttt{f\_range\_late}, and
\texttt{f\_range\_ratio}.}
The range (maximum minus minimum) of raw fitness values in each window, and
their ratio:
\begin{align}
\label{eq:frange}
  F_{\mathrm{E}} &= \max_{t \le T/4} y_t
                    - \min_{t \le T/4} y_t, \\
  F_{\mathrm{L}} &= \max_{t > 3T/4} y_t
                    - \min_{t > 3T/4} y_t, \\
  R_F            &= F_{\mathrm{L}} \;/\; F_{\mathrm{E}}.
\end{align}
Here $F_{\mathrm{E}}$ and $F_{\mathrm{L}}$ are the early and late
fitness ranges, and $R_F$ is their ratio.
A shrinking fitness range ($R_F < 1$) indicates the algorithm is
converging
to a narrow fitness band near the optimum. A persistent or growing range
suggests the algorithm still samples points of widely varying quality.

% ============================================================================
\section{Additional Trajectory Features}
\label{sec:features:novel}

The six features in this section adapt established statistical and time-series concepts to the analysis of metaheuristic search trajectories. In each case the underlying mathematical tool (autocorrelation, plateau detection, coefficient of variation, Pearson correlation, per-dimension range analysis) is standard in adjacent fields, but its application as a behavioural feature for metaheuristic characterisation is, to our knowledge, new to the optimisation literature.

\paragraph{\texttt{step\_size\_autocorrelation}.}
\label{sec:features:novel:ssac}

Autocorrelation of fitness values is a well-studied landscape descriptor
\cite{weinberger1990fitness}, and we include it above as
\texttt{fitness\_autocorrelation}.  The closest precedent for
autocorrelation as a behavioural feature in optimisation is de~Nobel et
al.~\cite{deNobel2021tsfresh}, who applied lag-1 autocorrelation to the
CMA-ES conjugate evolution path $\lVert p_\sigma \rVert$ as part of a
\texttt{tsfresh}-extracted feature set.  Our feature differs in two
ways.  First, it operates on Euclidean displacement step sizes from the
trajectory, which exist for any algorithm, rather than on a CMA-ES
internal control signal.  Second, it is selected and reported as a
single named, interpretable behavioural descriptor rather than embedded
in an automatically extracted feature bank.

Let $s_t = \lVert x_{t+1} - x_t \rVert$ denote the step size at evaluation~$t$.
The lag-1 autocorrelation of the step-size series is:
\begin{equation}
\label{eq:ssac}
  \rho_s
    = \frac{\sum_{t=1}^{T-2}(s_t - \bar{s})(s_{t+1} - \bar{s})}
           {\sum_{t=1}^{T-1}(s_t - \bar{s})^2}.
\end{equation}

The step-size autocorrelation $\rho_s$ captures \emph{search momentum}. High positive autocorrelation (${\sim}0.9$) means the algorithm takes similarly-sized steps consecutively, characteristic of smooth step-size adaptation mechanisms such as the $1/5$-th success rule or CMA-ES's path-length control. Low or negative autocorrelation indicates erratic step-size variation, typical of random restarts or poorly
tuned mutation operators.

\paragraph{\texttt{fitness\_plateau\_fraction}.}
\label{sec:features:novel:plateau}

Existing stagnation metrics in BLADE (\texttt{longest\_no\_improvement\_streak},
\texttt{last\_improvement\_fraction}) measure plateaus in the \emph{best-so-far}
curve.  This conflates two distinct phenomena, the algorithm may be sampling diverse points that happen to be worse (active exploration), or it may be stuck in a genuinely flat region of the landscape.  We disentangle these by measuring plateaus in the \emph{raw} fitness sequence.

Let $\epsilon = 10^{-8} \cdot (\max_t y_t - \min_t y_t)$ be a relative
tolerance.  The plateau fraction is:
\begin{equation}
\label{eq:plateau}
  P_{\mathrm{flat}}
    = \frac{1}{T-1}\sum_{t=1}^{T-1}
      \mathbf{1}\!\bigl[|y_{t+1} - y_t| < \epsilon\bigr].
\end{equation}

High plateau fraction~$P_{\mathrm{flat}}$ means the algorithm repeatedly evaluates points with nearly identical fitness, indicating it is trapped in or deliberately exploiting a flat region.  This differs fundamentally from best-so-far stagnation, an algorithm can have a long no-improvement streak while still sampling diverse fitness values (because all are worse than the incumbent).  Conversely, a high plateau fraction with a short no-improvement streak indicates the algorithm has found a flat basin and is refining within it.

\paragraph{\texttt{half\_convergence\_time}.}
\label{sec:features:novel:halftime}

While \texttt{average\_convergence\_rate} captures the geometric mean of
per-step convergence ratios, it says nothing about when during the budget the
bulk of improvement occurs.  Two algorithms can have the same convergence rate
but very different temporal profiles. One may improve rapidly in the first 10\% of evaluations and then plateau, while another may start slowly but achieve a late breakthrough.

The half-convergence time measures when 50\% of the total best-so-far
improvement is reached:
\begin{equation}
\label{eq:halftime}
  t_{1/2}
    = \frac{1}{T-1}\min\!\bigl\{t : f^*_t \le f^*_1
      - \tfrac{1}{2}(f^*_1 - f^*_T)\bigr\}.
\end{equation}

Low values of the half-convergence time~$t_{1/2}$
(${\sim}0.001\text{--}0.01$) indicate rapid early convergence
followed by slow refinement. High values (${\sim}0.5\text{--}1.0$) indicate
a slow start with late breakthroughs.  The feature returns~$1.0$ when no
improvement occurs.

\paragraph{\texttt{improvement\_spatial\_correlation}.}
\label{sec:features:novel:impspatial}

We are not aware of a prior feature that relates the \emph{distance} an algorithm travels to the \emph{magnitude} of the improvement it achieves.  We define the improvement
spatial correlation as the Pearson correlation between step size and
improvement magnitude, restricted to improving steps:
\begin{align}
\label{eq:impspatial}
  \mathcal{I} &= \{t : y_t < f^*_{t-1}\}, \nonumber \\
  r_{\mathrm{IS}}
    &= r_{\mathrm{Pearson}}\!\bigl(
         \{s_{t-1}\}_{t \in \mathcal{I}},\;
         \{f^*_{t-1} - y_t\}_{t \in \mathcal{I}}
       \bigr).
\end{align}

Positive values of the improvement spatial
correlation~$r_{\mathrm{IS}}$ indicate that large jumps yield large
improvements.
Negative values indicate that small, local steps produce the biggest gains.  Values near zero mean improvement
magnitude is independent of the distance travelled.  This feature separates
algorithms that improve through directed long-range moves from those that
improve through local refinement.

\paragraph{\texttt{improvement\_burstiness}.}
\label{sec:features:novel:burst}

Barab\'asi~\cite{barabasi2005bursts} showed that many natural processes exhibit bursty dynamics. Events tend to cluster in time rather than arriving at a steady rate.
We apply this idea to the temporal pattern of fitness improvements.

Let $\mathcal{I} = \{t_1, t_2, \ldots, t_k\}$ be the sorted indices of
improving evaluations, and $\tau_i = t_{i+1} - t_i$ the inter-improvement
intervals.  The improvement burstiness is the coefficient of variation of these
intervals:
\begin{equation}
\label{eq:burst}
  B_{\mathrm{imp}}
    = \frac{\mathrm{std}(\tau)}{\mathrm{mean}(\tau)}.
\end{equation}

The improvement burstiness~$B_{\mathrm{imp}}$ is a coefficient of
variation (CV).  Values near~1 correspond to Poisson-like random
improvements.  CV $\gg 1$ indicates bursty dynamics where improvements arrive
in clusters separated by long stagnation phases.  CV $\ll 1$ indicates
metronomic, evenly-spaced improvements.  This feature characterises the
temporal \emph{texture} of the search process. Does the algorithm make steady
incremental progress, or does it alternate between productive bursts and idle
periods?

\paragraph{\texttt{dimension\_convergence\_heterogeneity}.}
\label{sec:features:novel:dimhet}

Standard convergence metrics (spread ratio, step-size trend) aggregate across
all dimensions.  They cannot detect whether the algorithm converges uniformly
in all dimensions or locks in on some dimensions while continuing to explore
others, a pattern that reflects sensitivity to problem separability.

For each dimension~$j$, we compute the range shrinkage ratio between the early
and late windows:
\begin{equation}
\label{eq:shrinkage}
  \sigma_j
    = \frac{\max_{t > 3T/4} x_t^{(j)} - \min_{t > 3T/4} x_t^{(j)}}
           {\max_{t \le T/4} x_t^{(j)} - \min_{t \le T/4} x_t^{(j)}}.
\end{equation}
The convergence heterogeneity is the standard deviation of these ratios across
dimensions:
\begin{equation}
\label{eq:dimhet}
  H_{\mathrm{dim}}
    = \mathrm{std}\!\bigl(\sigma_0, \sigma_1,
      \ldots, \sigma_{d-1}\bigr).
\end{equation}

Low dimension convergence heterogeneity~$H_{\mathrm{dim}}$ means all
dimensions converge at a similar rate, so the
algorithm treats the problem as isotropic.  High heterogeneity means some
dimensions converge tightly while others remain spread, suggesting the algorithm
identifies and exploits coordinate-wise structure.

% ============================================================================
\section{Computational Cost Analysis}
\label{sec:features:cost}

\Cref{tab:feature_cost} summarises the computational cost of each feature
category, measured over 50 repeated runs on a synthetic trajectory matching
the experimental setup ($T = 10{,}000$ evaluations, $d = 5$).  Each candidate
algorithm is evaluated on 50 MA-BBOB trajectories (10 functions
$\times$ 5 instances), so per-candidate cost is $50\times$ the per-trajectory
cost.  The dominant cost is the existing BLADE metrics, which build a KD-tree
for dispersion and iterate over all evaluations for nearest-neighbour
distance.  The 17 features from categories 1, 4, and~5 add under 250\,ms per
candidate in total.

\begin{table}[t]
  \centering
  \caption{Computational cost of behavioural feature extraction.  Timings
    are the mean of 50 runs on a single CPU core for $T = 10{,}000$ evaluations
    in $d = 5$ dimensions.  SampEn uses subsampled trajectories ($N = 1{,}000$).}
  \label{tab:feature_cost}
  \footnotesize
  \begin{tabular}{@{}lrrl@{}}
    \toprule
    Category & \# & Per traj. (ms) & Per cand. \\
    \midrule
    Existing (BLADE)            & 11 & $743.7 \pm 5.3$   & $\sim$37\,s   \\
    Step-size dynamics          &  4 & $1.6 \pm 0.1$     & $\sim$80\,ms  \\
    Information-theoretic       &  4 & $27.5 \pm 0.5$    & $\sim$1.4\,s  \\
    Adapted pop.\ dynamics      &  7 & $1.0 \pm 0.1$     & $\sim$48\,ms  \\
    Novel features              &  6 & $1.5 \pm 0.1$     & $\sim$75\,ms  \\
    \midrule
    \textbf{Total}              & \textbf{32} & $\mathbf{775.2 \pm 5.4}$ & $\sim$\textbf{39\,s} \\
    \bottomrule
  \end{tabular}
\end{table}

Relative to the MA-BBOB evaluation itself (30--120\,s per candidate to run
the algorithm 50 times), the total feature extraction overhead of
${\sim}39$\,s per candidate is modest.  Over a full experimental run of 100
candidates, this amounts to roughly 65 minutes of feature extraction compared
to 1--3 hours of algorithm evaluation.
